%% 第二章 载人登月逃逸典型任务分析与剖面设计

\chapter{载人登月逃逸典型任务分析与剖面设计}
\chaptermark{载人登月逃逸分析与剖面设计}

\section{载人登月逃逸任务特殊性}

载人登月任务因为涉及月球轨道，和近地轨道（LEO）载人飞行任务相比有很大区别，这也让载人登月的逃逸剖面设计有了特殊性。表~\ref{tab:compare}~列出了载人登月与近地轨道逃逸任务的主要差异。

\begin{table}[htbp]
\caption{载人登月与近地轨道逃逸任务对比}
\label{tab:compare}
\begin{tabularx}{\textwidth}{P{3cm}CC}
\toprule
对比维度 & 近地轨道任务 & 载人登月任务 \\
\midrule
任务周期 & 数小时至数天 & 数天至数周 \\
飞行阶段 & 发射上升、在轨运行、返回再入 & 发射上升、地月转移、环月轨道、月面下降/上升、月地转移、再入 \\
最大飞行速度 & 约 $7.8\ \mathrm{km/s}$（轨道速度） & 约 $11.2\ \mathrm{km/s}$（第二宇宙速度） \\
引力环境 & 地球引力 & 地球+月球双引力场 \\
逃逸窗口 & 每一圈均有返回机会 & 地月相对位置变化周期长，窗口有限 \\
通信时延 & 低，近乎即时 & 长，单向延迟大约 $1.3\ \mathrm s$ \\
在轨救援 & 可快速发射飞船救援 & 救援周期长 \\
\bottomrule
\end{tabularx}
\end{table}

这些差异给载人登月逃逸任务带来了不少设计难点。

首先是全任务链覆盖。从发射台到月面再返回地球，登月任务的飞行阶段比近地轨道任务多很多。各阶段的动力学环境、推进系统状态和逃逸手段都有显著差异，每个阶段都需要设计独立的逃逸剖面。

地月转移段的双引力场耦合问题。飞船飞到地月转移段时，会同时受到地球和月球引力的作用，逃逸弹道设计必须同时考虑这两个天体的引力影响。单一引力源假设在这个阶段的误差太大，得用限制性三体模型来替代。

通信延迟问题。地月距离大概38万公里，信号单向时延能达到1.3秒左右。地面没法实时操控飞船或传达指令，必须让飞船自主做决策。这对飞船的自主故障检测和应急处理能力有较高要求。

返回窗口比较有限。近地轨道任务的飞船每90分钟就有一次返回机会；而地月转移和环月轨道段的返回窗口受地月相对位置制约，返回窗口间隔很长，能达到数小时甚至数天一次。所以逃逸方案需要具备长时间的轨道保持能力。

\section{发射阶段逃逸剖面}

发射阶段是逃逸系统使用概率最高的阶段，也是飞船动力学环境最严酷的阶段。按逃逸时的高度和飞行状态分类，可以把发射阶段逃逸分成三种典型剖面。

\subsection{待发段逃逸（Pad Abort）}

待发段逃逸发生在火箭还没起飞或刚离开发射台时，处于零高度、零速度状态。这时火箭可能出现推进剂泄漏、发射台火灾或爆炸等故障，逃逸系统必须在几秒内把乘员舱带离危险区域。

\begin{figure}[htbp]
\centering
\subfloat[中国载人飞船逃逸系统飞行试验]{
\label{fig:cm-se-pad}
\includegraphics[width=0.48\textwidth]{中国载人飞船逃逸系统零高度飞行试验瞬间.jpg}}
\hfill
\subfloat[猎户座飞船PA-1飞行试验]{
\label{fig:orion-pad}
\includegraphics[width=0.48\textwidth]{猎户座飞船零高度逃逸(PA-1)飞行试验实况.png}}
\caption{待发段逃逸飞行试验}
\label{fig:pad-abort-test}
\end{figure}

逃逸过程如下：故障触发后，逃逸主发动机和姿控发动机同时点火。AM提供约1.96MN的巨大推力，让组合体快速获得向上的加速度。大概2秒内，逃逸塔完成从垂直起飞到建立水平速度分量的俯仰转弯。ACM通过姿控喷气保持组合体姿态稳定，按预定程序分阶段完成俯仰转弯、下程保持和姿态重定向。组合体在约1.6km高度处完成抛塔发动机点火与LAS分离，之后乘员舱继续独立飞行上升到约1.8km最高点，再转入伞降段。从逃逸指令下达到抛塔的整个过程约13秒。

设计特点有三个：零初始速度条件下需要极高推重比，AM推力约为起飞重量的14倍；因为气动阻尼较弱，姿态控制主要靠ACM，控制系统的响应速度和精度很关键；除此之外，初始俯仰角接近90度（垂直状态），欧拉角运动学方程处于近奇异区，需要数值保护措施。

\subsection{最大动压区逃逸（Max-Q Abort）}

最大动压区是火箭上升过程中气动载荷最大的阶段，这时逃逸的主要困难在于高速、高动压和跨声速非线性气动效应~\cite{苗志文2026}。

\begin{figure}[htbp]
\centering
\includegraphics[width=0.85\textwidth]{猎户座飞船最大动压逃逸(AA-2)飞行状态示意图.png}
\caption{猎户座飞船最大动压逃逸(AA-2)飞行状态示意}
\label{fig:aa2-flight}
\end{figure}

逃逸过程：火箭飞行约55秒、高度约9.4km、马赫数约1.18时触发逃逸。AM和ACM同时点火，组合体与火箭分离后在推力作用下继续加速爬升。因为初始已有约50度的弹道倾角和明显的来流速度，逃逸后的气动环境非常复杂——跨声速区激波位置变化会导致气动系数剧烈波动，高动压（约30kPa）下气动恢复力矩给姿态控制提供了较大的裕度，但也让气动参数的不确定性直接转化为弹道偏差。经过大约15秒的逃逸拉离和10秒姿态重定向后，JM在27秒抛塔，乘员舱独立飞行。

设计特点包括：跨声速区（M≈1.0~1.2）气动系数随马赫数剧烈变化，CX和Cyα需要分段建模；高动压条件让气动效应对扰动高度敏感，气动参数偏差是影响弹道精度的最主要因素；另外，气动恢复力矩给姿态提供了自然阻尼，角速度变化幅度很小，姿态裕度充足。

\subsection{高空逃逸}

高空逃逸发生在逃逸塔抛离之后、飞船入轨之前。这个阶段飞行器已经处于较高的高度（通常>40km），大气非常稀薄，气动力作用可以忽略。逃逸塔完成使命，已经被抛离，逃逸需要靠飞船服务舱（SM）的主发动机，属于整船逃逸模式。

逃逸过程如下：在故障触发后，SM主发动机先点火，让飞船与故障火箭快速分离并拉开安全距离。随后飞船用姿控发动机调整姿态，选择进入合适的再入返回窗口。再入弹道类型有两种：一种是弹道式再入，另一种是跳跃式再入，弹道式再入的飞船直接以接近零升力的方式再入大气层，过载较大但轨迹简单可靠。跳跃式再入的飞船以较小的再入角进入大气层上边界，利用升力短暂跳出后再二次进入，可以有效降低峰值过载和热流，适合初始能量接近轨道速度的工况。

设计关键在于：这个阶段剩余的推进剂质量有限，必须精确计算机动所需的ΔV，避免推进剂不足导致无法到达预定的再入点。再入走廊的宽度通常只有1度到2度，再入角太大的话，峰值热流和过载会超出航天员生理极限和热防护上限；如若再入角太小的话，又可能弹出大气层，导致无法正常着陆。高空逃逸的弹道终点需要与搜救区域匹配，还得避开人口稠密的陆地区域。

\section{轨道飞行阶段逃逸剖面}

\subsection{地月转移轨道逃逸}

地月转移段飞行器处于地球与月球引力的过渡区域，只采用地球引力源假设误差太大，需要引入月球引力，用限制性三体模型来计算。

逃逸方案有两种典型的返回路径：一种是直接返回，用服务舱发动机反推减速，让飞行器脱离地月转移轨道，沿自由返回轨迹飞向地球；另一种是绕月返回，继续沿地月转移轨道飞到月球附近，借助月球引力弹弓效应改变飞行方向，进入返回地球的轨道。阿波罗13号事故用的就是绕月返回方案，把登月舱当作救生艇，最终成功返回。

这个阶段的设计难点是，返回轨道的计算必须同时考虑地球和月球引力，而且返回窗口受地月相对几何关系的约束。

\subsection{环月轨道逃逸}

环月轨道段飞行器已经进入月球引力主导的区域。这时如果出现推进系统故障或生命保障系统异常，飞船应该尽快脱离环月轨道返回地球。

逃逸方案的时序包括计算最佳返回窗口、启动轨道控制发动机调整轨道至月地转移轨道以及最终返回地球。

\section{登月阶段逃逸剖面}

登月阶段是整个登月任务中操作最复杂、风险最高的阶段之一，主要包括月面下降、月面停留和月面上升三个子阶段。

\subsection{月面下降段逃逸}

月面下降段按逃逸触发时机，可以分为悬停前逃逸和悬停后逃逸两种。

悬停前逃逸：当下降发动机推力异常、导航系统偏差或姿态控制失效时，紧急关闭下降发动机，马上启动上升发动机，以垂直向上姿态快速爬升。爬升到安全高度后调整姿态进入环月轨道，与指令舱交会对接。

悬停后逃逸：飞行器已经到达月面附近（悬停高度），这时逃逸面临更大的障碍物规避需求。优先延长悬停时间来快速规划规避轨迹；如果无法规避，就启动上升发动机，以小角度倾斜爬升避开障碍物，再转入环月轨道。

\subsection{月面停留段逃逸}

航天员在月面活动时，可能遇到生命保障系统失效、辐射超标等突发情况。这个阶段的逃逸方案要求快速撤离、简化流程，立即启动月面上升器脱离月面，最终进入环月轨道与指令舱汇合。

月面紧急起飞不需要像地球发射那样克服稠密大气的阻力，但月面重力只有地球的1/6，上升轨迹的设计要和月球重力适配，考虑低重力环境的动力学特性。上升器与指令舱的交会对接策略也需要提前规划，确保对接窗口能够匹配。

\subsection{月面上升段逃逸}

上升段如果遇到主发动机故障、姿态控制失效或导航偏差等情况，应该立即切换到备用推进系统或备用控制模式。逃逸剖面要能保证上升器进入安全的环月轨道，即要求即使主推力损失，也能通过延长燃烧时间或用备份发动机实现入轨。

\section{返回阶段逃逸剖面}

\subsection{月地转移段逃逸}

这个阶段飞行器已经进入返回地球的轨道，面临的逃逸场景主要有推进系统异常、导航偏差和热防护系统预警。和地月转移段类似，返回轨道修正需要借助地球引力场来计算，精确调整再入角，确保返回舱在安全的再入走廊内进入地球大气层。

\subsection{地球再入段逃逸}

再入段是任务的最后一个高风险阶段。返回舱会以接近第二宇宙速度（约11.2km/s）进入地球大气层，防热系统面临1640摄氏度以上的高温考验。如果再入角太大，峰值过载和热流可能超过航天员生理极限和防热能力极限；要是再入角太小，返回舱可能弹出大气层，导致无法正常着陆。

逃逸方案：如果防热盾损伤较轻，可以调整姿态减小再入角，降低峰值热流和过载，用跳跃式再入（Skip Re-entry）轨迹来延长减速过程；如果降落伞系统异常，应该切换到备用伞系，保证着陆速度在安全范围内。
